\documentclass[11pt]{article}
\usepackage[margin=1in]{geometry}
\usepackage{amsmath,amssymb,amsthm,array,booktabs,hyperref,microtype}
\hypersetup{colorlinks=true,urlcolor=blue}
\title{The cost of simultaneous torsion specialization\\over abelian local fields}
\author{A companion to the exponent-seven census}\date{}
\setlength{\emergencystretch}{3em}
\begin{document}\maketitle
\begin{abstract}
A finite family of integral polynomials is tested at primitive roots of unity translated by one. We determine the smallest combined valuation budget that can force excess above a prescribed tolerance at every point. The formula covers every finite abelian extension of $\mathbb Q_p$ and every finite set of orders. The dyadic case distinguishes odd and even polynomial degrees below the cyclotomic conductor.
\end{abstract}
\section{Exact torsion selection over local abelian fields}
\label{sec:quadratic-tower}
A local ingredient of the Lean proof specializes integral polynomials at
characters of prime-power order. The incoming report
\emph{Sharp torsion specialization, random symplectic blocks, and the reach
of one-sided Frobenius tests} computed an exact fourth-root table and asked for
higher-order formulas. We obtain the entire tower for each quadratic
extension of $\mathbb Q_2$. The proof also gives a descent principle valid
at every prime. These are local results and their proofs do not require RH.

Normalize $v(2)=1$ here. Write $N_h=2^{h-1}$ and $r_h=1/N_h$;
the test points on shell $h$ are the $z=\zeta-1$ of exact order
$2^h$ for $\zeta$. Fix a nonempty finite set $S$ of shell indices and
$\rho\ge\max_{h\in S}r_h$. A packet consists of nonzero
$g_i\in\mathcal O_K[T]$ and witnesses $x_i\in\mathbb C_2$ with
$v(x_i)\ge\rho$ and finite $t_i=v(g_i(x_i))$. Its cost is $\sum_i t_i$.
It blocks $S$ at tolerance $\delta\ge0$ when, at every test point,
$\sum_i\max(0,v(g_i(z))-t_i)>\delta$.
Let $C_K(S,\delta)$ be the least blocking cost.

\paragraph{Quadratic tower theorem.}
Let $K=\mathbb Q_2(\sqrt D)$ for a representative in the table below.
One has $C_K(S,\delta)=\sum_{h\in S}c_{D,h}(\delta)$, with
$c_{D,1}(\delta)=1$ for all $D,\delta$. On the other shells the cost is
\[
 c_{D,h}(\delta)=
 \begin{cases}1/2&0\le\delta<t_{D,h},\\1&\delta\ge t_{D,h}.
 \end{cases}
\qquad
\begin{array}{c|cc}
D&t_{D,2}&t_{D,h}\ (h\ge3)\\\hline
5&0&0\\
-1,\ 3&1/2&1/2\\
2,\ -2,\ 10,\ -10&1/4&1
\end{array}
\]
In particular a budget $B$ guarantees a common point with loss at most
$\delta$ exactly when $B<C_K(S,\delta)$. At equality a packet with every
witness equal to zero attains the obstruction. The fourth-root column is
the previous report's result. The stabilized column covers all higher shells.

\paragraph{Reduction to a monic polynomial on one shell.}
For general $p$ normalize $v(p)=1$, and write $r=1/\varphi(p^h)$ on
shell $h$. The Gauss valuation of $g=b\prod(T-\alpha)$ at weight $u$
is $G_u(g)=v(b)+\sum_\alpha\min(u,v(\alpha))$. For $u=r\le\rho$ and
$v(x_i)\ge\rho$ one has $t_i\ge G_\rho(g_i)\ge G_r(g_i)$.
Only roots with $v(\alpha)=r$ can raise $v(g_i(z))$ above $G_r(g_i)$.
Collect all such roots, with multiplicity, into a monic $P_h$ of degree
$m_h$. Their multiset is Galois-stable, so $P_h\in\mathcal O_K[T]$.
It follows that the loss on this shell is at most $v(P_h(z))-m_hr$.
Also $\sum_h m_hr_h\le\sum_i(G_\rho(g_i)-G_0(g_i))\le\sum_i t_i$.
Conversely any such single-shell polynomials blocking their own shells,
used separately with witnesses zero, give the sum of their costs.
Thus the minimum on one shell is the first $mr$ for which a monic
$P\in\mathcal O_K[T]$ of degree $m$ has
\[
                      v(P(z))>mr+\delta\quad\text{at every test point}.
\]
Indeed, removing all roots of valuation different from $r$ preserves this
inequality with the remaining degree, since each removed root contributes
at most $r$ to the evaluation. A first successful polynomial therefore
has all its roots on the radius, and $mr=v(P(0))\in v(K^\times)$.
The shifted cyclotomic polynomial always succeeds at cost one.
This reduction and the budget additivity are from the earlier report.

\paragraph{Tame descent of the entire degree test.}
Let $L/K$ be any finite tamely ramified extension of local fields of residue
characteristic $p$. Every successful monic degree-$m$ test polynomial over
$L$ gives one of the same degree over $K$ on any finite Galois-stable set of
test points with a common threshold. Choose $u\in\mathcal O_L$ with
$\operatorname{Tr}_{L/K}u=1$. Such a $u$ exists because the trace for
the maximal unramified subextension surjects integrally (reduce the trace
of its integers to the trace of finite fields), and the remaining total
ramification degree is a unit. Explicitly, divide a trace-one element
of the unramified subextension by that degree. Take the trace
coefficientwise of $uP$. The result is integral and monic. Each conjugate of
$P$ satisfies the strict inequality at every point by stability of the set.
The ultrametric inequality gives the same strict inequality for the trace.
Extension of scalars gives the reverse implication. In particular every
prime-power shell cost at every tolerance is invariant under tame base
change. The same argument works in several variables for a common
evaluation threshold and an imposed monic coefficient.

\paragraph{Quadratic reduction.}
The unramified field $U=\mathbb Q_2(\sqrt5)$ identifies each of
$\mathbb Q_2(\sqrt{-1})$, $\mathbb Q_2(\sqrt2)$ and $\mathbb Q_2(\sqrt{-2})$
with its unramified quadratic twist after adjoining $U$. For the first
pair, $3/(-5)$ is a square since it is $1$ modulo $8$ as a $2$-adic unit.
Tame descent equates costs within the three pairs in the table.
For $U$ the cost equals that over $\mathbb Q_2$, where $mr\in\mathbb Z$
forces it to be at least one. For any ramified quadratic field and $h\ge2$,
$mr\in\tfrac12\mathbb Z$, so the only candidate below one is $m=N_h/2$.
For $h=1$ the positive integral degree itself is the cost.

\paragraph{Extremizers and lower bounds on their losses.}
Set $X=1+T$. Over $\mathbb Q_2(i)$ take $X-i$ at $h=2$.
Its minimum evaluation valuation at the primitive fourth roots is one.
Its cost is $1/2$, so its minimum excess is $1/2$.
Over $\mathbb Q_2(\sqrt{\pm2})$, choose the linear center
$a=1+\sqrt{\pm2}$ on that shell. Since $v(a^2+1)=3/2$ and
$v(i-(-i))=1$, both $v(a-i)$ and $v(a+i)$ equal $3/4$.
No center in this quadratic field can improve the minimum: the difference
from $a$ of any other center has valuation in $\tfrac12\mathbb Z$ or is
zero, so neither distance can increase past $3/4$.
Over $\mathbb Q_2(i)$ the mutual distance of the two targets bounds the
minimum by one. This proves the fourth-root column.

For $h=3$, over $\mathbb Q_2(\sqrt2)$ and
$\mathbb Q_2(\sqrt{-2})$ respectively use
\[
                     X^2-\sqrt2 X+1,\qquad X^2-\sqrt{-2} X-1 .
\]
Choose the signs of the square roots compatibly with a primitive $\xi$ of
order eight. The two pairs of roots are $\{\xi,\xi^7\}$ and
$\{\xi,\xi^3\}$ respectively. At each omitted primitive root the two
distances have valuations $1$ and $1/2$. Hence the minimum evaluation
valuation is $3/2$, and the excess over cost $1/2$ is exactly one.
Any successful construction at level $s$ lifts to every $h\ge s$ by
replacing $X$ by $X^{2^{h-s}}$. The primitive points map surjectively,
the degree scales by $2^{h-s}$, and its product with $r_h$ stays fixed.
Every root of these lifted polynomials in the $T$ coordinate is on its
designated radius, since the base roots used in the lifts for
$s=2$ over $\mathbb Q_2(i)$ and $s=3$ over $\mathbb Q_2(\sqrt{\pm2})$
come from primitive roots of unity. This supplies all strict lower sides
of the thresholds in the table.

\paragraph{Upper bounds at every higher level.}
The following short orbit argument specializes the earlier report's
cyclotomic-coefficient descent. Write $F_s=\mathbb Q_2(\zeta_{2^s})$
and $h\ge s$. Suppose an eligible polynomial of degree $m=N_h/2$ over
a subfield of $F_s$ has all its roots $\alpha$ on the test radius.
Group roots into orbits over $F_s$, retaining multiplicities. An orbit of
degree $d$ costs $d/N_h$ in valuation, which lies in
$N_s^{-1}\mathbb Z$ by the norm. Give this orbit
$dN_s/N_h$ marks. Assign its marks to a unit class $b$ modulo $2^s$ by
choosing a nearest test point $\xi_h^b-1$ for one root. Each conjugate
has a nearest point with the same exponent class modulo $2^s$, since
automorphisms fixing $F_s$ act by exponents congruent to $1$ modulo $2^s$.

There are $N_s/2$ marks in total. For an unmarked class $a$, put
$k=v_2(a-b)<s$. At any lift of $a$ the orbit contributes at most
$(2^k-1)/N_s$ in excess per mark: the nearest-point choice bounds each
root distance valuation by $v(\xi_h^a-\xi_h^{b'})=2^k/N_h$,
where $b'$ is its own nearest exponent. Subtract the baseline $1/N_h$
per root. At $s=2$ there is one mark, and the other class has $k=1$,
giving excess at most $1/2$. At $s=3$ there are two marks. If their
classes agree modulo four, choose the opposite class modulo four;
both $k$'s equal one. Otherwise choose an unmarked lift of either class
modulo four, giving $k$'s two and one. The excess is therefore at most
$(3+1)/4=1$. Since $\mathbb Q_2(\sqrt{\pm2})\subset F_3$,
these bounds apply to both claimed stabilized thresholds. At equality the
strict blocking inequality fails, completing the proof.

\paragraph{Example.}
For $K=\mathbb Q_2(\sqrt{10})$ and $\delta=3/4$, the fourth-root
shell costs one and every higher shell costs $1/2$. The shell set
$\{2,3,4,5\}$ costs $5/2$ exactly. Over $\mathbb Q_2(\sqrt3)$ the
same four shells instead cost four. At $\delta=1$ every shell costs
one in every quadratic field. The file \texttt{eighth\_roots.py}
independently recovers the first stabilized column by finite arithmetic.

\subsection*{Abelian coefficient fields}
Tame descent, together with the cyclotomic calculation in
\emph{Torsion specialization and sparse prime sums} (the incoming
archive \texttt{grh-torsion-and-prime-sampling.zip}), gives a general
positive-tolerance formula. For any prime put $N_j=(p-1)p^{j-1}$ and
\[
 S_j(L)=\frac1{N_j}\sum_{i=1}^{j-1}N_i
          \left\lfloor L/N_i\right\rfloor\quad(0\le L<N_j),
 \qquad S_j(N_j)=+\infty .
\]
Let $m_j(\delta)$ be the first $L\in\{0,\ldots,N_j\}$ with
$S_j(L)>\delta$. The cited cyclotomic descent theorem gives the cost over
$F_s=\mathbb Q_p(\mu_{p^s})$ on shell $h$ as
$m_j(\delta)/N_j$, where $j=\min(h,s)$. More precisely, when $h\ge s$
and the degree spent on that shell has normalized cost $L/N_s$, the
minimal excess somewhere on the shell is at most $S_s(L)$. The bound
comes from selecting successively a least-occupied unit residue class
and then a least-occupied lift; the number of marks in its class
modulo $p^i$ is at most $\lfloor L/N_i\rfloor$. The orbit construction
in the previous proof gives the same compression of roots to marks at
level $s$. Balanced markings attain the bound.

\paragraph{All odd primes.}
For odd $p$, every finite abelian $K/\mathbb Q_p$, every eligible finite
$S$ and every real $\delta\ge0$, write $a=v_p(e(K/\mathbb Q_p))$.
Then
\[
 C_K(S,\delta)=\sum_{h\in S}
   \frac{m_{\min(h,a+1)}(\delta)}{N_{\min(h,a+1)}} .
\]
In particular the entire tolerance profile depends only on the wild part
of the ramification index.
By the local Kronecker--Weber theorem take $K\subset F_n U$ with
$U/\mathbb Q_p$ finite unramified and $n\ge a+1$.
The groups of the two linearly disjoint factors identify
$\operatorname{Gal}(F_n U/U)$ with $(\mathbb Z/p^n)^\times$, which
is cyclic. Moreover $K\cap U$ is the maximal unramified subfield of $K$,
so $[KU:U]=e(K/\mathbb Q_p)$. This degree divides $(p-1)p^a$.
Uniqueness of each intermediate degree in a cyclic extension gives
$KU\subset F_{a+1}U$. The extension $F_{a+1}U/KU$ has degree prime to
$p$, and $KU/K$ and $F_{a+1}U/F_{a+1}$ are unramified. Tame descent
therefore identifies all the costs with those for $F_{a+1}$.
For the field-theoretic input see the proof of local Kronecker--Weber in
A. Sutherland, \href{https://math.mit.edu/classes/18.785/2017fa/LectureNotes20.pdf}
{\emph{Number Theory I}, Lecture 20}.

\paragraph{All abelian dyadic fields on high shells.}
By the same compositum argument, for $p=2$ there is a subfield
$A\subset F_s$ with $KU=AU$ after an unramified base change. Choose
$s\ge1$ minimal with $A\subset F_s$. When $A=F_s$ the preceding
cyclotomic formula applies on every shell. In the remaining case $s\ge3$,
and for every $h\ge s$ the exact cost is
\[
 \frac1{N_s}\min\{L\in\{0,2,4,\ldots,N_s\}: S_s(L)>\delta\}.
\]
Indeed $H=\operatorname{Gal}(F_s/A)$ avoids the nonidentity kernel of
reduction from level $s$ to $s-1$. The unit group modulo $2^s$
is $C_2\times C_{2^{s-2}}$. Every element of order at least four has a
power equal to $1+2^{s-1}$. Of the four elements of order dividing two,
a subgroup avoiding $1+2^{s-1}$ has order at most two. Hence $H$ has
order two, generated by $c=-1$ or $c=-1+2^{s-1}$ modulo $2^s$.
The possible costs $L/N_s$ have $L$ even, by
$e(A/\mathbb Q_2)=N_s/2$ and the constant term of a shell polynomial.
The cyclotomic bound over $F_s$ gives the stated lower bound for the cost.

Conversely for even $L\le N_s$ select the $L/2$ exponents
$1+4i$ ($0\le i<L/2$) and their images under multiplication by $c$.
These are distinct: $c\equiv3\pmod4$.
The resulting marks are $H$-stable and each unit class modulo $2^j$
contains at least $\lfloor L/N_j\rfloor$ marks. For $j\ge2$, the first
half cycles through all classes congruent to one modulo four with period
$2^{j-2}$; the second half cycles through the complementary classes with
the same period. At level one every mark lies in the only class.
The monic polynomial with roots $\xi_s^b-1$ at the selected marks
is therefore defined over $A$. An unmarked point of exponent $u$ has
excess given by summing $(2^{v_2(u-b)}-1)/N_s$ over the marks.
The lower occupancy bounds give excess at least $S_s(L)$; at marked
points it is infinite. Pullback by $X\mapsto X^{2^{h-s}}$ in the
$X=1+T$ variable realizes the same cost at all higher levels.
Tame descent transfers this conclusion from $A$ to $K$.
The remaining lower shells also admit a formula, as follows.

\paragraph{Dyadic shells below the coefficient conductor.}
Continue with $A$ of index two, and let $2\le h<s$. Define the
following extremal loss profile on $0\le m\le N_h$:
\[
 E_h(m)=
 \begin{cases}
 S_h(m),&m\text{ even},\\
 S_h(m-1)+r_h-r_s,&m\text{ odd}.
 \end{cases}
\]
Then the exact cost on shell $h$ is $m/N_h$ for the first $m$ with
$E_h(m)>\delta$. Shell one always costs one. Together with the high-shell
formula and the odd-prime theorem this evaluates every shell for every
finite abelian coefficient field over $\mathbb Q_p$.

For even $m$ the unrestricted upper bound is $S_h(m)$. At level $h$,
$c$ acts on exponents by negation. The construction with $m/2$ marks
in the class one modulo four and their negatives therefore realizes
the bound over $A$. For odd degree the obstruction is the best
approximation to a primitive target by an element of odd degree over $A$.
Fix $t=\xi_s-1$ and $z=\xi_h-1$, using primitive roots at those levels.
In the quadratic extension $F_s/A$ write $z=u+bt$, $u,b\in A$. One has
\[
 v(z-c(z))=2r_h,\qquad v(t-c(t))=2r_s,\qquad
 v(z-u)=v(bt)=2r_h-r_s=: \lambda .
\]
Indeed $v_2(c-1)=1$, and $c(z)$, $c(t)$ denote the Galois images.
The difference of any two distinct test points on shell $h$ has valuation
at least $2r_h$. Thus every test point has distance valuation $\lambda$
from $u$. Since $h<s$ gives $\lambda>r_h$, this single center is
integral, of valuation $r_h$, and attains excess $r_h-r_s$ everywhere.

Moreover $\lambda$ is an odd integral multiple of $r_s$. If $[A(\alpha):A]$
is odd, the ramification index $e'$ of that extension is odd. Its value
group $(2r_s/e')\mathbb Z$ does not contain $\lambda$.
Consequently $v(\alpha-u)\ne\lambda$, and the ultrametric inequality
bounds $v(\alpha-z')\le\lambda$ at every test point $z'$.
A monic polynomial over $A$ of odd degree has an irreducible factor of
odd degree. Take one of its roots $\alpha$ and apply the unrestricted
bound $S_h(m-1)$ to the remaining roots. This gives the claimed upper
bound for its minimum excess. Conversely multiply the extremizing
even-degree polynomial by $T-u$. Its minimum excess gains exactly
$r_h-r_s$. The constructed polynomials have every root on the designated
radius and hence attain their costs with witness zero.
Tame descent gives the same profile of feasible monic degree tests
after adjoining $U$ and over the original field $K$.

\paragraph{Checks and a sharp order bound.}
Two supplied scripts check the dyadic constructions.
The first validates 1,036 balanced configurations and exhausts the stable
subsets for levels through five. The second independently solves linear
congruences over the dyadic integers. At conductors three through six it
checks 84 degree, shell and involution choices, verifying each optimum is
attained and the next valuation step is impossible. It represents
$\mathcal O_{F_s}$ in the basis of powers of $t=\xi_s-1$; powers of
the uniformizer $t\,c(t)$ give an integral basis of $A$.
As an odd-prime example, if $p=3$, $v_3(e(K))=1$ and $K/\mathbb Q_3$
is abelian, every shell $h\ge2$ costs $2/3$ at tolerance $1/3$.
Indeed $N_2=6$ and $S_2(L)=\lfloor L/2\rfloor/3$ for $L<6$,
so the strict threshold first occurs at $L=4$. Four consecutive shells
therefore guarantee a successful specialization for a packet of budget
at most two; three can all fail at exactly that budget.

\paragraph{Sources.} The Lean input that motivates the question is
\href{https://github.com/lsernau/grh-lean/blob/72fead102f663d5c7a407964a13dc0cf9d977207/grh/src/New/Congruences/Selmer/RootsOfUnity.lean}{\texttt{RootsOfUnity.lean}}.
The two predecessors used in the proofs are
\href{https://github.com/VladimirReshetnikov/ProveIt/blob/d31039a48d963692e5a5176f47b6603fc9e9d9cf/docs/incoming/grh-torsion-and-prime-sampling.zip}{the cyclotomic-field calculation}
and \href{https://github.com/VladimirReshetnikov/ProveIt/blob/d31039a48d963692e5a5176f47b6603fc9e9d9cf/docs/incoming/grh_sharp_torsion_and_symplectic_blocks.zip}{the fourth-root table and fixed-field monic test}.
\end{document}
